% 原书分章PDF第11—12页：第27.4节，章末。
\section{注记与延伸阅读}

本章内容基于 Díaz 和 van der Laan（2013b）以及 Luedtke 等（2015b）先前的工作；希望了解更多讨论和细节的读者，可参阅这些文献。我们强调，领域专家应当仅根据背景科学知识，就能判断敏感性参数 $\delta$ 的每个取值是否合理。我们也强调，必须使用统计上有效的方法估计 $\psi_0$。特别是，我们建议采用靶向学习。在这一步施加缺乏依据的参数假设，是一种常见做法；它通常会增加统计偏差，从而使敏感性分析的结论失效。

下面讨论与本章方法相关的一些方法。Rosenbaum 和 Rubin（1983a）考虑了平均处理效应参数模型的估计，其假定是：对某个未测量混杂因素作调整后，该参数便可识别。VanderWeele 和 Arah（2011）在相同的未测量混杂因素假设下，给出了用统计参数表示因果参数时的偏差的一般公式，从而改进了这些结果。这些方法有时需要很强的简化假设，例如假定这一未测量混杂因素是二分类的；在这些假设下，可以得到可解释的敏感性参数。为了获得因果参数的不确定性界，需要先获得对已测量混杂因素作调整的统计参数的置信界，因此，需要采用现代统计方法，才能保证这些置信界有效。类似方法也曾用于得到交互作用参数的界（VanderWeele 等，2012a）与直接效应的界（VanderWeele，2010）。Ding 和 VanderWeele（2015）中的界正是通过这一思路定义的；当未测量混杂因素并非二分类时，它所需的参数远少于早期方法。

这些敏感性分析程序与部分识别文献有关。后者在很弱的假设下，甚至不作假设时，推导因果参数与统计参数之差的界，例如只利用结局的界（Manski，1990，2003；Horowitz 和 Manski，2000；MacLehose 等，2005）。这些分析得到的界一旦能提供有用信息，便很有说服力；但在许多情况下，它们可能过于保守，甚至不能提供效应符号方面的信息。在例27.2的背景下，Manski（2003）详细讨论了可归结为利用结局已知界的因果参数界。该文献还给出了其他假设下的界，例如单调性假设；这些假设可能有助于启发条件 $\delta$ 的其他选择。

Horowitz 和 Manski（2000）使用了类似于第27.3节的方法，在部分可识别性问题中，为这些界获得不确定性区域（uncertainty region）。Woutersen（2006）考虑了在部分可识别性问题中，给定上下界的渐近线性估计量后，如何构造这样的不确定性区域。实际上，这两项工作都考虑了一种改进程序，它保证以趋于 $0.95$ 的概率覆盖参数 $\psi_{\mathrm{causal}}$，而不是覆盖已知包含 $\psi_{\mathrm{causal}}$ 的整个区域 $(\psi_0+L_{\mathrm{causal}}(\delta),\psi_0+U_{\mathrm{causal}}(\delta))$。敏感性分析也可作类似改进，但这里不作探讨。我们在第27.3节提出的不确定性区域，在思路上类似于 Vansteelandt 等（2006）给出的区域，但不需要为识别偏差设定一个缺乏合理性的半参数模型。
